\chapter{Implementación de los patrones en Flink SQL}
\label{anx:patrones}

Cada patrón se organiza en los mismos pasos, cada uno en un fichero SQL que se ejecuta con el
cliente de Flink y una definición común de tablas y vistas que se carga al inicio:

\begin{enumerate}
    \item \textbf{Flujo}: publica en un \textit{topic} intermedio
    (\texttt{det.p1\_flujo}, \texttt{det.p2\_flujo}, \texttt{det.p3\_flujo}) la evidencia del
    patrón, sin repeticiones y en orden de tiempo de evento.
    \item \textbf{Campañas}: correlación en dos niveles sobre ese flujo.
    \item \textbf{Emisión}: construcción del hallazgo OCSF y publicación en
    \texttt{det.campaigns}.
    \item \textbf{Flujo continuo}: el mismo patrón en modo \textit{streaming}, con una ventana
    de sesión como segundo nivel, frente a su resultado por lotes.
\end{enumerate}

El Listado~\ref{lst:p1} reproduce los dos niveles de P1 en su versión por lotes.

\begin{lstlisting}[language=SQL,caption={P1: transiciones (nivel 1) y campañas (nivel 2).},label={lst:p1}]
-- NIVEL 1: dos evidencias consecutivas del mismo equipo a <= 10 min,
-- al menos una de credenciales (transicion bidireccional).
CREATE TEMPORARY VIEW transiciones_p1 AS
SELECT *
FROM (SELECT dataset, host, ts, `time` AS t_ms, cred FROM p1_flujo)
MATCH_RECOGNIZE (
    PARTITION BY dataset, host
    ORDER BY ts
    MEASURES E.t_ms AS t_ini_ms, C.t_ms AS t_fin_ms
    ONE ROW PER MATCH
    AFTER MATCH SKIP TO NEXT ROW
    PATTERN (E C)
    DEFINE
        C AS (C.cred = 1 OR E.cred = 1)
         AND TIMESTAMPDIFF(SECOND, E.ts, C.ts) <= 600
) AS t;

-- NIVEL 2: un hueco de mas de 10 min entre transiciones abre campana.
CREATE TEMPORARY VIEW transiciones_marcadas_p1 AS
SELECT *,
       CASE WHEN LAG(t_fin_ms) OVER (PARTITION BY dataset, host
                                      ORDER BY t_ini_ms) IS NULL
              OR (t_ini_ms - LAG(t_fin_ms) OVER (PARTITION BY dataset,
                                      host ORDER BY t_ini_ms)) > 600000
            THEN 1 ELSE 0 END AS abre
FROM transiciones_p1;

CREATE TEMPORARY VIEW campanas_base_p1 AS
SELECT dataset, host, campana,
       MIN(t_ini_ms) AS inicio, MAX(t_fin_ms) AS fin,
       COUNT(*) AS transiciones
FROM (SELECT *, SUM(abre) OVER (PARTITION BY dataset, host
             ORDER BY t_ini_ms ROWS BETWEEN UNBOUNDED PRECEDING
             AND CURRENT ROW) AS campana
      FROM transiciones_marcadas_p1)
GROUP BY dataset, host, campana;
\end{lstlisting}

En modo continuo, el segundo nivel se sustituye por una ventana de sesión sobre el tiempo de
cada transición, que Flink proporciona como \texttt{MATCH\_ROWTIME()}:

\begin{lstlisting}[language=SQL,caption={P3: segundo nivel en modo continuo.},label={lst:p3-sesion}]
SELECT dataset, host,
       MIN(t_ini_ms) AS inicio, MAX(t_fin_ms) AS fin,
       COUNT(*) AS transiciones
FROM TABLE(SESSION(TABLE transiciones_p3_rt PARTITION BY clave,
                   DESCRIPTOR(rt), INTERVAL '10' MINUTE))
GROUP BY clave, dataset, host, window_start, window_end
HAVING SUM(con_otra) >= 1;
\end{lstlisting}

\section{Restricciones de Flink 2.2 encontradas}

\begin{itemize}
    \item No se admiten cuantificadores voraces al final de un patrón, lo que obliga a la
    correlación en dos niveles (Sección~\ref{sec:dos-niveles}).
    \item \texttt{DEFINE} no admite subconsultas y \texttt{MEASURES} no admite
    \texttt{COUNT(DISTINCT)}.
    \item \texttt{WITHIN} limita la duración total de la coincidencia, no el hueco entre
    elementos.
    \item \texttt{LAG} no acepta marcas de tiempo con zona; se opera sobre milisegundos.
    \item Una ventana de sesión sobre una vista de una tabla cuyo tiempo es una columna
    calculada pierde el atributo de tiempo; se aplica sobre la tabla.
    \item Una vista que contiene una ventana de sesión con partición de varias columnas no se
    puede volver a expandir al consultarla; se escribe como subconsulta.
    \item \texttt{JSON\_OBJECT} inserta como texto los objetos que llegan de una subconsulta;
    el documento de campaña se compone con \texttt{LISTAGG} y \texttt{CONCAT}.
    \item Los literales de texto en un \texttt{CASE} se rellenan con espacios hasta la mayor
    longitud; se convierten explícitamente a \texttt{STRING}.
\end{itemize}
